\documentclass[10pt,letterpaper,twocolumn]{article}

% ===================== PAQUETES =====================
\usepackage[spanish,es-tabla]{babel}
\usepackage[utf8]{inputenc}
\usepackage[T1]{fontenc}
\usepackage{lmodern}
\usepackage{microtype}
\usepackage{geometry}
\usepackage{amsmath,amssymb,amsfonts,mathtools}
\usepackage{graphicx}
\usepackage{booktabs}
\usepackage{array}
\usepackage{subcaption}
\usepackage{multirow}
\usepackage{enumitem}
\usepackage{xcolor}
\usepackage{float}
\usepackage{caption}
\usepackage{subcaption}
\usepackage{subcaption}
\usepackage{hyperref}
\usepackage[most]{tcolorbox}
\usepackage{fancyhdr}
\usepackage{titlesec}
\usepackage{setspace}
\usepackage{csquotes}
\usepackage{tabularx}
\usepackage{makecell}

% ===================== CONFIGURACION =====================
\geometry{left=1.65cm,right=1.65cm,top=2.0cm,bottom=2.0cm,columnsep=0.72cm}
\setlength{\parskip}{0.42em}
\setlength{\parindent}{0pt}
\setlength{\headheight}{14pt}
\singlespacing

\definecolor{azulpinn}{HTML}{12355B}
\definecolor{azulclaro}{HTML}{EAF2FA}
\definecolor{gris}{HTML}{F4F4F4}
\definecolor{rojo}{HTML}{B23A48}
\definecolor{verde}{HTML}{287D4F}
\definecolor{amarillo}{HTML}{FFF4CC}

\hypersetup{colorlinks=true,linkcolor=azulpinn,citecolor=azulpinn,urlcolor=azulpinn}

\pagestyle{fancy}
\fancyhf{}
\lhead{Identificación de ecuaciones para flujos de fluidos mediante redes neuronales
informadas por la física}
\rhead{Proyecto Final}
\cfoot{\thepage}


% ===================== COMANDOS =====================
\newcommand{\placeholderfigure}[2]{%
\begin{figure}[H]
\centering
\begin{tcolorbox}[colback=gris,colframe=azulpinn,width=0.96\columnwidth,height=#2,arc=2mm]
\centering
\vspace*{0.05cm}
\textbf{ESPACIO PARA FIGURA}\par
\vspace{0.1cm}
\small #1
\end{tcolorbox}
\caption{#1}
\end{figure}
}

\newcommand{\placeholdertable}[2]{%
\begin{table}[H]
\centering
\begin{tcolorbox}[colback=gris,colframe=azulpinn,width=0.96\columnwidth,height=#2,arc=2mm]
\centering
\vspace*{0.05cm}
\textbf{ESPACIO PARA TABLA}\par
\vspace{0.1cm}
\small #1
\end{tcolorbox}
\caption{#1}
\end{table}
}

\newtcolorbox{notebox}{colback=azulclaro,colframe=azulpinn,arc=2mm,boxrule=0.8pt}
\newtcolorbox{warningbox}{colback=amarillo,colframe=rojo,arc=2mm,boxrule=0.8pt}

% ===================== DATOS DEL DOCUMENTO =====================
\title{\textbf{Identificación de parámetros en la ecuación de Burgers 2D mediante redes neuronales informadas por la física}\\

\author{Juan Jose Salazar Bedoya \and Cristian Camilo Victoria Parra}}
\date{16 de junio de 2026}

\begin{document}
\maketitle

\begin{abstract}
Este documento organiza el informe final a partir del articulo \textit{Equation Identification for Fluid Flows via Physics-Informed Neural Networks} y del notebook de reproducción desarrollado en Google Colab. El problema estudiado es inverso: se parte de datos parciales de una solución de la ecuación vectorial de Burgers en dos dimensiones y se busca reconstruir el campo de velocidades $\mathbf{u}(x,y,t)=(u,v)$, ademas de estimar los parámetros físicos de conveccion $\alpha$ y difusión $\nu$. El paper genera soluciones de referencia con un método clásico de volúmenes finitos sobre una malla de $256\times 256$ puntos espaciales y $72$ tiempos, mientras que el notebook permite ejecutar tres modos: \texttt{smoke}, \texttt{demo} y \texttt{paper}. La versión ejecutada en Colab usa el modo \texttt{demo}, con malla reducida $48\times48\times24$, arquitectura menor y menos pasos de entrenamiento, por lo que sus resultados deben interpretarse como una reproducción didáctica y no como una replica exacta de escala completa. El informe compara explícitamente los datos del paper con los datos usados en el notebook, deja espacios para insertar las figuras y tablas exportadas desde Colab, y estructura la discusión alrededor de los tres controles exigidos por el proyecto: comparación con referencia clásica, sensibilidad y plausibilidad física.
\end{abstract}

\textbf{Palabras clave:} PINNs, Burgers 2D, problema inverso, identificación de parámetros, SciML, volúmenes finitos.

% ============================================================
\section{Introducción}

Las Physics-Informed Neural Networks (PINNs) son una clase de redes neuronales diseñadas para problemas descritos mediante ecuaciones diferenciales y se toman en cuenta las leyes físicas (leyes de newton, termodinámica, ecuaciones de Maxwell...) el modelo en cuestión no solo toma en cuenta el aprendizaje habitual, si no que también toma en cuenta que esos datos sean físicamente coherentes. La motivación para emplear PINNs surge de la necesidad de obtener grandes volúmenes de datos de calidad, así como problemas en los que la física es bien conocida pero la solución analítica es difícil de abordar, haciendo necesario su uso.\\

El artículo estudia un problema muy interesante relacionado con el comportamiento de los fluidos. En lugar de asumir que todos los parámetros físicos son conocidos y simplemente calcular la solución de la ecuación, los autores consideran la situación contraria: se cuenta únicamente con algunas observaciones del sistema y, a partir de ellas, se intenta determinar qué parámetros dieron origen a ese comportamiento. En decir, el objetivo no es solo describir el fenómeno, sino también descubrir información sobre el modelo físico que lo gobierna.\\

Como banco de pruebas, el artículo utiliza la ecuación de Burgers en dos dimensiones. Aunque es más simple que las ecuaciones completas que describen el movimiento de los fluidos, conserva características importantes de estos fenómenos, como el transporte y la difusión. Esto la convierte en una buena herramienta para estudiar el comportamiento.

\subsection{Objetivo general}

Principalmente se quiere replicar de manera parcial la metodología y avances del paper para identificar parámetros en la ecuación de Burgers 2D mediante PINNs. Se quieren analizar los resultados generados en la réplica, compararlos con los del notebook y los aportados por el método clásico.

\subsection{Objetivos específicos}

\begin{itemize}[leftmargin=*]
    \item Presentar el problema físico asociado a la ecuación de Burgers 2D.
    \item Formular las ecuaciones diferenciales parciales (EDPs), condiciones iniciales y condiciones de frontera usadas en el paper.
    \item Comparar los datos, mallas, hiperparámetros y tamaños de entrenamiento del paper con los usados en el notebook.
    \item Describir la formulación PINN para reconstruir $\mathbf{u}(x,y,t)$ y estimar $\alpha$ y $\nu$.
    \item Insertar y analizar las figuras y tablas generadas en Colab.
    \item Evaluar los resultados mediante comparación con referencia, sensibilidad y plausibilidad física.
\end{itemize}

% ============================================================
\section{Problema físico y contexto del paper}

Los autores trabajan con ecuaciones que describen como cambia una variable física en el espacio y en el tiempo. Se puede pensar en:

\begin{itemize}
    \item Velocidad de un fluido
    \item Temperatura de un material
    \item Concentración de una sustancia
    \item Propagación de una onda
\end{itemize}

En todos estos casos se presenta un cambio en el espacio y tiempo.
\\
Para este caso, el problema físico correspondo a la evolución de un campo de velocidades bidimensional dentro del cuadrado unitario, el cual, para entenderlo, se puede imaginar un cuadro, no uno dibujado en una hoja, sino una región del espacio donde hay un fluido. En cada punto de ese cuadro se coloca una flecha, donde:

\begin{itemize}[leftmargin=*]
    \item La dirección de la flecha dice hacia dónde se mueve el fluido
    \item La longitud de la flecha indica qué tan rápido se mueve el fluido
\end{itemize}

Visto, por ejemplo, de esta forma:

\[
\begin{matrix}
\nearrow & \rightarrow & \rightarrow \\
\uparrow & \nearrow & \rightarrow \\
\uparrow & \uparrow & \nearrow
\end{matrix}
\]

Ahora viene una palabra importante: \textbf{evolución}. Ese mismo recuadro puede cambiar a:
\[
\begin{matrix}
\rightarrow & \rightarrow & \searrow \\
\nearrow & \rightarrow & \searrow \\
\uparrow & \nearrow & \rightarrow
\end{matrix}
\]

Entonces el problema ya no  es sobre cómo es el campo de velocidades, sino cómo cambia el campo de velocidades conforme pase el tiempo, es decir, la evolución temporal de un campo de velocidades bidimensional. Hablamos, además, de un cuadro unitario porque los autores decidieron estudiar el fenómeno únicamente en la región $[0 ,  1]\times[0 , 1]$.

La solución es un campo vectorial
\[
\mathbf{u}(x,y,t)=(u(x,y,t),v(x,y,t)),
\]
donde $u$ es la componente horizontal y $v$ la componente vertical de la velocidad.\\

La evolución del campo de velocidades está gobernada por la ecuación de Burgers bidimensional. Esta ecuación combina dos mecanismos físicos importantes: la convección y la difusión. La convección puede entenderse como el transporte de información por el propio flujo, de forma similar a cómo una corriente arrastra una gota de tinta. La difusión, por su parte, representa los efectos viscosos y tiende a suavizar las variaciones del campo.\\

Desde el punto de vista del paper, el reto no es generar la solución, pues esa solución se puede obtener por un método clásico de volúmenes finitos. El reto es usar solo una fracción de esos datos para recuperar los parámetros de la ecuación. Por eso el problema es inverso: se conocen algunos valores de $\mathbf{u}$, pero no necesariamente se conocen $\alpha$ y $\nu$.

\begin{figure}[H]
\centering
\begin{subfigure}[t]{0.48\linewidth}
    \centering
    \includegraphics[width=\linewidth]{figura1-paper.png}
    \caption{Figura del paper.}
    \label{fig:figura1-paper}
\end{subfigure}
\hfill
\begin{subfigure}[t]{0.48\linewidth}
    \centering
    \includegraphics[width=\linewidth]{figura1-notebook.png}
    \caption{Figura del notebook.}
    \label{fig:figura1-notebook}
\end{subfigure}

\vspace{0.15cm}

{\small
Instantáneas del componente $u(x,y,t)$ para $\alpha = 2$ y $\nu = 10^{-4}$ en varios tiempos.
}
\caption{Comparación entre la Figura 1 del paper y la reproducción generada en Colab.}
\label{fig:comparacion-figura1}

\end{figure}
En la Figura 1 se comparan las instantáneas temporales de la componente $u(x,y,t)$ dadas en el paper con las generadas en el notebook. En ambos casos se observa que, a medida que avanza el tiempo, el patrón inicial deja de variar simplemente en la dirección $y$, sino que se empiezan a formar estructuras asociadas al flujo rotacional del campo vectorial $\mathbf{u}=(u,v)$. Esta característica es importante porque el paper no busca estudiar un caso trivial de la ecuación de Burgers en 2D, sino un lugar donde la dinámica presenta rotacional no nulo, por lo cual la identificación de parámetros resulta más compleja.\\

En la reproducción del notebook, las estructuras rotacionales son menos evidentes que en la figura del paper. Esto puede explicarse por la reducción en la escala computacional, ya que se trabaja con una malla más pequeña, menos pasos temporales y una configuración de entrenamiento más ligera. Por esta razón, la figura del notebook debe verse como una interpretación didáctica y cualitativa, no como una replica exacta.

% ============================================================
\section{Ecuaciones gobernantes}

\subsection{Ecuación vectorial de Burgers 2D}

La ecuación considerada en el paper es
\begin{equation}
    \mathbf{u}_t+\alpha(\mathbf{u}\cdot\nabla)\mathbf{u}-\nu\Delta\mathbf{u}=0,
    \qquad \mathbf{u}=(u,v).
    \label{eq:burgers_vectorial}
\end{equation}
Aquí aparecen exactamente las tres cosas que se han venido mencionando: 
$$\text{Cambio temporal} + \text{Convección} + \text{Difusión} = 0.$$

Es decir, $\alpha$ controla el peso del termino convectivo y $\nu$ controla la difusión viscosa. El operador convectivo es
\begin{equation}
    (\mathbf{u}\cdot\nabla)=u\frac{\partial}{\partial x}+v\frac{\partial}{\partial y},
\end{equation}

Por ejemplo, si un objeto se mueve en una montaña, la altura del terreno cambia, si se camina hacia el norte se siente una pendiente, pero si se camina hacia el este, el objeto notará otra otra pendiente. El operador convectivo responde qué pendiente sentirá según la dirección en la que se mueve el objeto.\\

Ahora, el operador laplaciano espacial es
\begin{equation}
    \Delta = \frac{\partial^2}{\partial x^2}+\frac{\partial^2}{\partial y^2}.
\end{equation}

Si se quisiese entender este operador con la analogía de la montaña sería de la siguiente forma; si estamos parados en un punto pronunciado de la montaña, todo a nuestro alrededor está más abajo. La naturaleza suele desgastar este pico con el tiempo.

\begin{itemize}[leftmargin=*]
    \item La lluvia 
    \item El viento
    \item La erosión
\end{itemize}

El pico se va aplanando.\\

Ahora, si se considera la contraparte, un valle profundo, todo lo demás está más arriba y de igual manera, con el tiempo, tierra que llegue sobre el valle lo rellenará para estar a la par. El operador laplaciano mide precisamente esto. qué tan diferente es un punto respecto a sus vecinos, si es muy alto o muy bajo en comparación a los demás, el operador querrá reducir esa diferencia.

Para lo cual, para este caso no se habla de altura, se habla de velocidades.

En componentes, el sistema queda
\begin{align}
    u_t+\alpha(uu_x+vu_y)-\nu(u_{xx}+u_{yy})&=0,\label{eq:fu}\\
    v_t+\alpha(uv_x+vv_y)-\nu(v_{xx}+v_{yy})&=0.\label{eq:fv}
\end{align}

\subsection{Dominio y condiciones}

El dominio espacio-temporal del paper es
\begin{equation}
    (x,y,t)\in [0,1]\times[0,1]\times[0,0.5].
\end{equation}

Las condiciones iniciales son
\begin{align}
    u(x,y,0)&=\sin(6\pi y)x(1-x),\label{eq:ic_u}\\
    v(x,y,0)&=-\sin(6\pi x)y(1-y).\label{eq:ic_v}
\end{align}

Totalmente útiles, pues la ecuación de burguers sólo nos dicen cómo evoluciona el sistema, no cómo empieza.

Las condiciones de frontera son homogéneas:
\begin{equation}
    \mathbf{u}=0\qquad \text{en }\partial([0,1]^2).
    \label{eq:bc}
\end{equation}

Estas condiciones son importantes porque se anulan en los bordes del cuadrado, pero producen un flujo rotacional no trivial en el interior. Esto evita que el problema sea demasiado simple y permite estudiar un régimen donde la identificación de parámetros no es inmediata, es decir, diferentes direcciones en el interior. 

% ============================================================
\section{Datos usados: paper frente a notebook}

En esta parte se verán las diferencias entre los datos usados en el paper y los datos usados en el notebook. A pesar de que la formulación matemáticas sea la misma, los resultados serán diferentes, ya que, como se dijo anteriormente, se tuvieron que trabajar con datos menores para reducir el costo computacional y temporal.

\subsection{Parámetros físicos del paper}

El paper usa diez configuraciones de parámetros. La viscosidad $\nu$ cambia por ordenes de magnitud y $\alpha$ cambia linealmente. La razón $\alpha/\nu$ sirve como indicador de dificultad: valores grandes corresponden a espacios mas convectivos y menos viscosos.

\begin{table}[H]
\centering
\caption{Configuraciones de parámetros del paper y del notebook.}
\begin{tabular}{ccc}
\toprule
$\nu$ & $\alpha$ & $\alpha/\nu$ \\
\midrule
$10^{-4}$ & 1.0 & 10000 \\
$10^{-4}$ & 1.5 & 15000 \\
$10^{-4}$ & 2.0 & 20000 \\
$10^{-3}$ & 1.0 & 1000 \\
$10^{-3}$ & 1.5 & 1500 \\
$10^{-3}$ & 2.0 & 2000 \\
$10^{-2}$ & 0.5 & 50 \\
$10^{-2}$ & 1.0 & 100 \\
$10^{-2}$ & 1.5 & 150 \\
$10^{-2}$ & 2.0 & 200 \\
\bottomrule
\end{tabular}
\end{table}

\subsection{Datos de referencia}
La escala utilizada en el notebook no es la misma que la escala usada en el paper. Esta reducción se hizo para reducir el costo computacional y el tiempo de ejecución en Colab. En particular, se disminuyeron la resolución espacial, el número de pasos temporales y algunos hiperparámetros de entrenamiento. Por lo tanto, el notebook reproduce la metodología general del artículo (generación de datos de referencia, muestreo de observaciones, entrenamiento de modelos PINN y comparación de errores), pero en una escala menor. Esto produjo que los resultados deben interpretarse como una forma de ver la estructura y como funciona lo planteado en el paper, y no como una replica exacta.
\begin{table}[H]
\centering
\caption{Comparacion entre la escala del paper y la escala ejecutada en el notebook.}
\begin{tabularx}{\columnwidth}{lXX}
\toprule
Elemento & Paper & Notebook \texttt{demo} \\
\midrule
Malla espacial & $256\times256$ & $48\times48$ \\
Tiempos & 72 & 24 \\
$t_{final}$ & 0.5 & 0.5 \\
Puntos por solucion & 4,718,592 & 55,296 \\
Barridos no lineales & reportado con solucion FiPy & 2 \\
Metodo de referencia & Volumenes finitos en FiPy & FiPy reducido \\
Objetivo & Benchmark completo & Reproduccion didactica \\
\bottomrule
\end{tabularx}
\end{table}

Esta diferencia de escala es importante para interpretar los resultados: si los errores del notebook son mayores que los del paper, esto no necesariamente indica que la metodología usada es la incorrecta, sino que es algo que se esperaría por trabajar con una escala menor a la utilizada en el paper.

\subsection{Tamaños de datos observados}

El paper separa dos experimentos: una PINN inversa y una linea base guiada por datos. La diferencia es que la PINN usa la EDP durante el entrenamiento, mientras que el modelo guiado por datos ajusta primero observaciones y despues estima los parametros mediante el residuo.

\begin{table}[H]
\centering
\caption{Tamaños de datos usados en el paper y en el notebook.}
\begin{tabularx}{\columnwidth}{lXX}
\toprule
Experimento & Paper & Notebook \texttt{demo} \\
\midrule
PINN inversa & 2048, 8192, 32768 & 512, 2048 \\
Data-driven & 128, 512, 2048, 8192, 32768, 131072 & 128, 512, 2048 \\
Configuraciones $(\alpha,\nu)$ & 10 & 10 \\
Optimizadores EDP & Adam y Newton & Adam y Newton \\
Total PINNs & $3\cdot2\cdot10=60$ & $2\cdot2\cdot10=40$ \\
Total data-driven & $6\cdot10=60$ & $3\cdot10=30$ \\
\bottomrule
\end{tabularx}
\end{table}

\subsection{Hiperparámetros del paper frente al notebook}

\begin{table}[H]
\centering
\caption{Hiperparámetros principales de la PINN.}
\begin{tabularx}{\columnwidth}{lXX}
\toprule
Elemento & Paper & Notebook \texttt{demo} \\
\midrule
Ancho de red & 256 & 80 \\
Profundidad & 10 & 4 \\
Activación & tanh & tanh \\
Pasos PINN & 100000 & 1500 \\
Batch IC & 1024 & 256 \\
Batch PDE & 2048 & 512 \\
$\lambda_{PDE}$ & 1 & 1 \\
$\lambda_{Data}$ & 10 & 10 \\
$\lambda_{IC}$ & 10 & 10 \\
Optimizador red & Adam & Adam \\
Optimizador PDE & Adam/Newton & Adam/Newton \\
LR inicial & $5\cdot10^{-3}$ & $3\cdot10^{-3}$ \\
LR mínimo & $10^{-5}$ & $10^{-5}$ \\
Puntos evaluación & no todos, muestra amplia & 6500 \\
\bottomrule
\end{tabularx}
\end{table}

\begin{table}[H]
\centering
\caption{Hiperparámetros principales del método guiado por datos.}
\begin{tabularx}{\columnwidth}{lXX}
\toprule
Elemento & Paper & Notebook \texttt{demo} \\
\midrule
Ancho de red & 256 & 80 \\
Profundidad & 10 & 4 \\
Activación & tanh & tanh \\
Épocas/pasos & 50000 & 1500 \\
Batch datos & 2048 & 512 \\
Batch IC & 2048 & 256 \\
Batch PDE & 8192 & 512 \\
$\lambda_{Data}$ & 1 & 1 \\
$\lambda_{IC}$ & 1 & 1 \\
Optimizador & Adam & Adam \\
LR inicial & $5\cdot10^{-3}$ & $3\cdot10^{-3}$ \\
LR mínimo & $10^{-6}$ & $10^{-5}$ \\
\bottomrule
\end{tabularx}
\end{table}
La diferencia en la escala justifica las diferencias numéricas entre el notebook y el paper, pero el resultado arrojado, a pesar de las diferencias debería mantener la tendencia mostrada en el paper, es decir, en los casos con menor viscosidad $\nu$, o equivalentemente con mayor valor $\alpha/\nu$, sean más complicados de aproximar, ya que la dinámica presenta estructuras más bruscas y el entrenamiento de la PINN se vuelve más exigente. También se espera observar que la regularización física, al incorporar el residuo de la EDP en la función de pérdida, aporta información adicional para la recuperación de los parámetros $\alpha$ y $\nu$, en comparación con un ajuste que se guía solo con datos.

% ============================================================
\section{Formulación PINN}

\subsection{Aproximación neuronal}

La red neuronal recibe como entrada el punto espacio-temporal $(x,y,t)$ y devuelve las dos componentes de la solución (u,v):
\begin{equation}
    (x,y,t)\longmapsto \mathbf{u}_\theta(x,y,t)=(u_\theta(x,y,t),v_\theta(x,y,t)).
\end{equation}

Para imponer las condiciones de frontera homogéneas de forma exacta, el paper usa una construcción con frontera dura:
\begin{align}
    u_\theta(x,y,t)&=x(1-x)y(1-y)NN_u(x,y,t;\theta),\\
    v_\theta(x,y,t)&=x(1-x)y(1-y)NN_v(x,y,t;\theta).
\end{align}

Así, cuando $x=0$, $x=1$, $y=0$ o $y=1$, ambas componentes se anulan automáticamente. Siendo esto una ventaja muy grande, pues así la red ya cumple las condiciones de frontera automáticamente. No hay que enseñárselas ni incluirlas en la función de pérdida.

\subsection{Residuo de la PDE}

La red de burguers impone una relación que debe cumplir la solución, el residuo nos dice cuánto se equivoca al red al intentar satisfacer esa ecuación.

El residuo para la primera componente es
\begin{equation}
\begin{split}
    f_u={}&u_{\theta,t}+\alpha\left(u_\theta u_{\theta,x}+v_\theta u_{\theta,y}\right)\\
    &-\nu\left(u_{\theta,xx}+u_{\theta,yy}\right),
\end{split}
\end{equation}
y el residuo para la segunda componente es
\begin{equation}
\begin{split}
    f_v={}&v_{\theta,t}+\alpha\left(u_\theta v_{\theta,x}+v_\theta v_{\theta,y}\right)\\
    &-\nu\left(v_{\theta,xx}+v_{\theta,yy}\right).
\end{split}
\end{equation}

Estas derivadas se calculan por diferenciación automática. El objetivo físico es que $f_u$ y $f_v$ sean pequeños en los puntos de colocación, es decir, una mayor aproximación de la red.

\subsection{Función de perdida}

La función de pérdida es un número que resume qué tan mal está haciendo su trabajo la red. Claramente queremos que ese te número sea lo más pequeño posible. 
La pérdida total usado por la PINN la coforman tres partes: datos, condición inicial y el residuo de la EDP.
\begin{equation}
    \mathcal{L}=\lambda_{PDE}\mathcal{L}_{PDE}+\lambda_{IC}\mathcal{L}_{IC}+\lambda_{Data}\mathcal{L}_{Data}.
\end{equation}

El termino de EDP es
\begin{equation}
    \mathcal{L}_{PDE}=\frac{1}{N_f}\sum_{i=1}^{N_f}\left(|f_u(x_i,y_i,t_i)|^2+|f_v(x_i,y_i,t_i)|^2\right).
\end{equation}

El termino de condición inicial es
\begin{equation}
\begin{split}
    \mathcal{L}_{IC}=\frac{1}{N_i}\sum_{i=1}^{N_i}&\left(|u_\theta(x_i,y_i,0)-u_0(x_i,y_i)|^2\right.\\
    &\left.+|v_\theta(x_i,y_i,0)-v_0(x_i,y_i)|^2\right).
\end{split}
\end{equation}

El termino de datos es
\begin{equation}
    \mathcal{L}_{Data}=\frac{1}{N_d}\sum_{i=1}^{N_d}\|\mathbf{u}_\theta(x_i,y_i,t_i)-\mathbf{u}(x_i,y_i,t_i)\|^2.
\end{equation}

El objetivo del problema inverso es encontrar los pesos de la red \(\theta\) y los parámetros físicos \(\alpha\) y \textit{v} (convección y viscocidad.
Esto se logra minimizando una función de pérdida \(\mathcal{L}\)
\begin{equation}
    \theta^*,\alpha^*,\nu^*=\arg\min_{\theta,\alpha,\nu}\mathcal{L}(\theta;\alpha,\nu).
\end{equation}
El propósito es encontrar los valores \(\theta\), \(\alpha\) y \textit{v} que hagan que la pérdida sea lo más pequeña posible.

El notebook no copia las fórmulas abstractas del paper, sino que las convierte en ecuaciones concretas que se pueden programar: usa $f_u$ , $f_v$ para la PDE, elimina las condiciones de frontera (porque van por separado con frontera dura), y separa la condición inicial como un término propio $\mathcal{L}_{IC}$ para reflejar mejor la implementación computacional.

% ============================================================
\section{Metodología computacional}

\subsection{Estructura del notebook}

El notebook esta organizado de la siguiente manera:
\begin{enumerate}[leftmargin=*]
    \item instalación de dependencias y configuración global;
    \item definición de los modos de ejecución \texttt{smoke}, \texttt{demo} y \texttt{paper};
    \item carga de la grilla de parámetros del paper;
    \item generación de datos con FiPy;
    \item muestreo de datos observados;
    \item definición de la red ED-MLP con frontera dura;
    \item calculo de derivadas automáticas y residuos;
    \item entrenamiento data-driven y entrenamiento PINN inverso;
    \item generación de figuras 1--8 y tablas comparativas.
\end{enumerate}

\subsection{Métricas de evaluación}

Para comparar la solución de referencia $\mathbf{u}$ con la predicción $\mathbf{u}_\theta$, se usa el error relativo
\begin{equation}
    E_{rel}=\frac{\left(\sum_{x\in X}\|\mathbf{u}(x)-\mathbf{u}_\theta(x)\|^2\right)^{1/2}}{\left(\sum_{x\in X}\|\mathbf{u}(x)\|^2\right)^{1/2}}.
\end{equation}

Para $\alpha$, el error es
\begin{equation}
    E_\alpha=\frac{|\widehat{\alpha}-\alpha|}{\alpha}.
\end{equation}

Para $\nu$; la viscosidad puede ser \( 10^{-2} \), \( 10^{-3} \) o \( 10^{-4} \). Son números en escalas muy diferentes. Si se usa la misma fórmula que para \( \alpha \), los errores no se entienden bien, por ello, en vez de comparar los números directamente, comparamos exponentes y hablamos de órdenes de magnitud. 

\begin{equation}
    E_\nu=\frac{|\log_{10}(\widehat{\nu}+\varepsilon)-\log_{10}(\nu+\varepsilon)|}{|\log_{10}(\nu+\varepsilon)|},
    \qquad \varepsilon=10^{-8}.
\end{equation}

Se busca un error menor al $10\%$, ya que es el umbral utilizado en el paper. Este se marca con una linea negra punteada en las gráficas.

% ============================================================
\section{Resultados del notebook y comparación con el paper}

\subsection{Resultado individual de prueba}

Antes de ejecutar todas las combinaciones de parámetros, el notebook hace una \textbf{prueba rápida} para verificar que todo funcione correctamente. Esta prueba usa:

\begin{itemize}
    \item \textbf{Convección} \( \alpha = 2.0 \)
    \item \textbf{Viscosidad} \( \nu = 10^{-4} \) (muy baja, fluido casi sin fricción)
    \item \textbf{Datos de entrenamiento}: 256 puntos 
    \item \textbf{Modo}: demo (no busca alta precisión)
\end{itemize}

El resultado obtenido fue:

\begin{table}[H]
\centering
\caption{Corrida individual de prueba registrada en el notebook.}
\begin{tabular}{lc}
\toprule
Variable & Valor \\
\midrule
Método & PINN \\
Optimizador PDE & Newton \\
$N_{data}$ & 256 \\
$\alpha$ real & 2.0 \\
$\nu$ real & $10^{-4}$ \\
$\widehat{\alpha}$ & 1.1372 \\
$\widehat{\nu}$ & 0.02854 \\
Error $\alpha$ & 0.4314 \\
Error $\nu$ & 0.6139 \\
Error relativo solución & 0.7798 \\
$\alpha/\nu$ & 20000 \\
\bottomrule
\end{tabular}
\end{table}

Esta prueba no busca igualar los resultados del paper; solo sirve para verificar que el código funciona, que no hay errores de programación. El error es alto porque se trata de una configuración pequeña, con pocos pasos y en el régimen mas difícil: $\alpha/\nu=20000$.

\begin{figure}[H]
    \centering
    \includegraphics[width=1\linewidth]{figura-perdida.png}
    \caption{Curva de entrenamiento de la corrida individual: perdida total, perdida EDP, perdida de datos y perdida de condición inicial.}
    \label{fig:placeholder}
\end{figure}

\subsection{Tabla 2: paper frente a reproducción disponible}

El paper reporta los siguientes promedios para las PINNs sobre el espacio de parametros:

\begin{table}[H]
\centering
\caption{Tabla 2 publicada en el paper.}
\begin{tabular}{lrrrr}
\toprule
Opt. PDE & Datos & Error $\alpha$ & Error sol. & Error $\nu$ \\
\midrule
Adam & 2048 & 0.080 & 0.152 & 0.093 \\
Newton & 2048 & 0.086 & 0.150 & 0.056 \\
Adam & 8192 & 0.085 & 0.139 & 0.097 \\
Newton & 8192 & 0.085 & 0.130 & 0.065 \\
Adam & 32768 & 0.067 & 0.139 & 0.101 \\
Newton & 32768 & 0.088 & 0.126 & 0.071 \\
\bottomrule
\end{tabular}
\end{table}

Con los resultados disponibles del notebook en modo \texttt{demo}, la tabla promedio obtenida fue:

\begin{table}[H]
\centering
\caption{Tabla 2 reproducida desde el notebook en modo \texttt{demo}.}
\begin{tabular}{lrrrr}
\toprule
Opt. EDP & Datos & Error $\alpha$ & Error sol. & Error $\nu$ \\
\midrule
Adam & 512 & 1.060 & 0.385 & 0.338 \\
Newton & 512 & 0.814 & 0.548 & 0.343 \\
Adam & 2048 & 1.135 & 0.980 & 0.767 \\
Newton & 2048 & 7.300 & 0.875 & 0.553 \\
\bottomrule
\end{tabular}
\end{table}

La comparación directa solo es posible para $N_{data}=2048$, pues es el único tamaño común entre la tabla publicada y la corrida \texttt{demo}. La diferencia fue:

\begin{table}[H]
\centering
\caption{Comparacion automatica entre paper y notebook para $N_{data}=2048$.}
\begin{tabular}{lrrrrrr}
\toprule
Opt. & $E_\alpha^p$ & $E_\alpha^n$ & $E_{sol}^p$ & $E_{sol}^n$ & $E_\nu^p$ & $E_\nu^n$ \\
\midrule
Adam & 0.080 & 1.135 & 0.152 & 0.980 & 0.093 & 0.767 \\
Newton & 0.086 & 7.300 & 0.150 & 0.875 & 0.056 & 0.553 \\
\bottomrule
\end{tabular}
\end{table}
La comparación muestra que los errores obtenidos en el notebook son considerablemente mayores que los reportados en el paper. Esta diferencia debe interpretarse teniendo en cuenta que la ejecución del notebook fue reducida: se usó una malla de $48\times48$, una red de ancho 80 y profundidad 4, y solo 1500 pasos de entrenamiento, mientras que el paper trabaja con una red de ancho 256, profundidad 10 y $100000$ pasos. Por tanto, la tabla no constituye una réplica numérica exacta del benchmark, sino una reproducción de menor escala de su metodología. La diferencia observada muestra el costo de reducir la resolución, la capacidad de la red y el número de iteraciones de entrenamiento.\\

En particular, el error de $\alpha$ con Newton aumenta de una gran manera en la reproducción del notebook. Esto sugiere que, bajo una configuración reducida, el método de estimación de parámetros puede volverse sensible a la calidad de la aproximación neuronal y al número limitado de pasos de entrenamiento.
\subsection{Figuras generadas en el notebook}
\begin{figure}[H]
    \centering
    \includegraphics[width=1\linewidth]{figura2.png}
    \caption{comparación PINN Newton frente a data-driven. Corresponde a la figura 2 del paper.}
    \label{fig:placeholder}
\end{figure}
Viendo la figura, se observa que el error relativo, el error para $\alpha$ y el error para $\nu$ es menor para \textit{método de Data-bases fitting}, lo cual difiere de los resultados del paper, ya que en el paper, el error generado es menor en las PINNs, es decir, para este caso, el resultado difiere del resultado del paper. Además, solo se trabaja con 2048 datos debido a que no se busca una replica exacta debido al costo computacional.
\begin{figure}[H]
    \centering
    \includegraphics[width=1\linewidth]{figura3.png}
    \caption{errores del método data-driven para diferentes cantidades de datos. Corresponde a la figura 3 del paper}
    \label{fig:placeholder}
\end{figure}
El error relativo es menor con valores de $\alpha$ bajos y valores de $\nu$ alto. Ahora, al estimar $\alpha$, los errores son muy similares, y son más alto del umbral buscado. Por otro lado, el error de $\nu$ baja mucho cuando $\alpha$ está entre $0.5$ y $1.5$ siempre y cuando el valor dado sea $0.001$ y $0.01$. Eso sí, cuando $\alpha$ es grande y $\nu$ vale $0.001$, se obtiene un menor porcentaje de error.

\begin{figure}[H]
    \centering
    \includegraphics[width=1\linewidth]{figura4.png}
    \caption{errores de parametros para PINNs con Adam y Newton. Corresponde a la figura 4 del paper}
    \label{fig:placeholder}
\end{figure}
Revisando la figura, los errores más bajos, tanto para $\alpha$ como para $\nu$ se obtienen con Newton y no con el optimizador de Adam, además, con Newton, $\alpha$ alcanzó el umbral de éxito con $\alpha=1.5;2$ $\nu=0.01;0.0001$. Los resultados difieren un poco con los del paper, ya que en el paper tuvieron errores muy similares.
\begin{figure}[H]
    \centering
    \includegraphics[width=1\linewidth]{figura5.png}
    \caption{error relativo de solucion para PINNs. Corresponde a la figura 5 del paper}
    \label{fig:placeholder}
\end{figure}
El error relativo para Newton y Adam es muy similar, tal como se muestra en el paper, con la diferencia de que en este caso, el error es demasiado alto.
\begin{figure}[H]
    \centering
    \includegraphics[width=1\linewidth]{figura6.png}
\end{figure}
\begin{figure}[H]
    \centering
    \includegraphics[width=1\linewidth]{figura7.png}
\end{figure}
\begin{figure}[H]
    \centering
    \includegraphics[width=1\linewidth]{figura8.png}
    \caption{Fresultados completos del baseline data-driven. Corresponden a las figuras 6, 7 y 8 del paper.}
\end{figure}

\subsection{Resultados obtenidos}

Los resultados obtenidos muestran que el notebook logra reproducir la estructura general del experimento del paper, pero no es igual de preciso. Esto es lógico, ya que la ejecución se hizo en modo demo, con una malla más pequeña, una red menos profunda y un número mucho menor de pasos de entrenamiento.\\

En la prueba, con $\alpha=2.0$, $\nu=10^{-4}$ y $N_{data}=256$, se obtuvo un error relativo de solución de $0.7798$, un error de $\alpha$ de $0.4314$ y un error de $\nu$ de $0.6139$. Estos valores son altos y están por encima del umbral del $10\%$ usado en el paper. Por lo tanto, esta prueba no muestra una buena estimación de los parámetros, sino que sirve para ver que el código funcione de buena manera.\\

En cuanto a la reconstrucción de la solución, los errores obtenidos en el notebook son mayores que los dados en el artículo. Para $N_{data}=2048$, el error relativo de solución fue $0.980$ con Adam y $0.875$ con Newton, mientras que en el paper los errores correspondientes son cercanos a $0.15$. Esto indica que la PINN en modo demo no logra reconstruir el campo de velocidades con la misma precisión que la configuración completa del paper. Aun así, lo planteado permite observar cómo se plantea el problema y cómo se comparan los métodos.\\

La estimación del parámetro $\alpha$ fue una de las partes más difíciles en la estimación. En todos los casos, el error promedio de $\alpha$ quedó por encima del $10\%$. En particular, para $N_{data}=2048$, el error fue $1.135$ con Adam y $7.300$ con Newton. Este resultado muestra que, con una configuración reducida, la estimación del coeficiente convectivo puede volverse muy sensible a la red y al número de pasos del entrenamiento. Si la red no aproxima bien la solución y sus derivadas, la estimación de los parámetros también se ve afectada.\\

Para el parámetro $\nu$, los errores también fueron mayores que los del paper, aunque en algunos casos se observa un comportamiento más favorable con Newton. Para $N_{data}=2048$, el error de $\nu$ fue $0.767$ con Adam y $0.553$ con Newton. Esto coincide parcialmente con la tendencia del artículo, donde Newton ayuda a mejorar la estimación del coeficiente difusivo. Sin embargo, los valores obtenidos siguen siendo altos, por lo que no se puede afirmar que $\nu$ haya sido recuperado de manera precisa.\\

Las figuras generadas en el notebook también muestran que los casos con mayor razón $\alpha/\nu$ tienden a ser más difíciles. Esto tiene sentido desde el punto de vista físico, porque cuando $\nu$ es pequeño la difusión suaviza menos la solución y el efecto convectivo domina más. En esos regímenes aparecen cambios más bruscos en el campo de velocidades, lo cual hace más difícil que la red aproxime correctamente la solución y recupere los parámetros.\\

En la comparación entre PINN y el método data-driven, los resultados del notebook no coinciden completamente con los del paper. En el artículo, las PINNs suelen tener ventaja para recuperar parámetros porque incorporan la ecuación diferencial en la función de pérdida. En la reproducción demo, en cambio, el método data-driven presenta en algunas gráficas errores menores. Esta diferencia no necesariamente contradice el paper, sino que muestra que la reducción de escala afecta bastante el comportamiento del experimento.\\

En resumen, los resultados obtenidos deben interpretarse como una reproducción didáctica y no como una réplica exacta del benchmark. El notebook reproduce las partes principales del procedimiento: generación de datos con FiPy, muestreo de observaciones, entrenamiento de la PINN, entrenamiento del método data-driven y comparación de errores. Sin embargo, la configuración demo no fue suficiente para alcanzar la precisión del paper. La principal conclusión es que la identificación de parámetros en la ecuación de Burgers 2D requiere una escala computacional mayor, especialmente cuando la viscosidad es baja y el cociente $\alpha/\nu$ es grande.
% ============================================================
\section{Extensión}
En esta parte, la extensión consiste en ver qué pasa cuando se reduce la escala computacional usada en el paper. Es decir, no se cambia la ecuación ni se propone una red nueva, sino que se analiza cómo cambian los resultados cuando el mismo método se ejecuta con una malla más pequeña, una red más sencilla y menos pasos de entrenamiento.\\

El paper trabaja con una configuración mucho más grande. En cambio, el notebook usa el modo demo, lo cual permite correr el experimento en Colab sin que el tiempo de ejecución sea tan alto. Sin embargo, esta reducción también afecta los resultados. Aunque la metodología es la misma, la red tiene menos capacidad para aprender la solución y para estimar correctamente los parámetros $\alpha$ y $\nu$.\\

Esto se puede ver en las tablas de resultados. Los errores obtenidos en el notebook son mucho mayores que los del paper, especialmente cuando se comparan los casos con $N_{data}=2048$. Por ejemplo, en el paper los errores son cercanos al umbral del $10\%$, mientras que en el notebook los errores quedan bastante por encima de ese valor. Esto muestra que reducir la escala ayuda a ejecutar el experimento, pero también hace que la identificación de parámetros sea menos precisa.\\

La estimación de $\alpha$ fue una de las más afectadas. En algunos casos, especialmente con Newton, el error aumenta demasiado. Esto indica que, si la red todavía no aprendió bien la solución, entonces las derivadas usadas para calcular el residuo de la EDP tampoco serán suficientemente buenas. Como los parámetros se estiman a partir de ese residuo, el error en la solución termina afectando también el error en $\alpha$ y $\nu$.\\

Por lo tanto, el costo computaciones influye mucho en este tipo de problemas, es decir, no basta con programar bien y entender la ecuación de Burgers, sino que se necesita una red lo suficientemente grande, una buena cantidad de puntos y un buen entrenamiento. El notebook (en modo demo), permite ver el funcionamiento y ver la idea general, pero los errores son altos. A continuación se muestra una comparación entre los datos usado en el paper y los usados en el modo demo del notebook:
\begin{table}[H]
\centering
\scriptsize
\setlength{\tabcolsep}{3pt}
\caption{Comparación entre la escala del paper y la escala ejecutada en el notebook.}
\begin{tabular}{lll}
\toprule
Elemento & Paper & Notebook demo \\
\midrule
Malla espacial & $256\times256$ & $48\times48$ \\
Tiempos & 72 & 24 \\
$t_{\mathrm{final}}$ & 0.5 & 0.5 \\
Puntos por solución & 4,718,592 & 55,296 \\
Método de referencia & Volúmenes finitos & FiPy reducido \\
Objetivo & Benchmark completo & Reproducción didáctica \\
\bottomrule
\end{tabular}
\label{tab:escala-paper-notebook}
\end{table}
Lo que se planteó se puede evidenciar en las tablas 6, 7, 8 y 9. También en las figuras de la sección 7.3.
\section{Validación obligatoria}

El proyecto no pide revisar solamente si la pérdida baja durante el entrenamiento. También es necesario mirar si la solución obtenida tiene sentido. Para esto se usan tres criterios: comparación con una referencia clásica, sensibilidad y plausibilidad física.

\subsection{Comparación con referencia clásica}

El primer criterio consiste en comparar los resultados de la PINN con una solución de referencia. En este caso, la referencia viene de FiPy, que resuelve la ecuación de Burgers usando un método clásico de volúmenes finitos. Esto es importante porque una pérdida pequeña no garantiza que la red haya aprendido correctamente la solución.\\

En el notebook, esta comparación se hace mediante el error relativo de solución. Este error mide qué tan lejos está la solución predicha por la red respecto a la solución generada con FiPy. Según los resultados obtenidos, los errores de solución en modo demo son altos. Por ejemplo, para $N_{data}=2048$, el error relativo de solución fue $0.980$ con Adam y $0.875$ con Newton. Estos valores son mayores que los reportados en el paper.

\begin{table}[H]
\centering
\scriptsize
\setlength{\tabcolsep}{6pt}
\caption{Comparación entre los errores de solución del paper y del notebook para $N_{data}=2048$.}
\begin{tabular}{lcc}
\toprule
Optimizador EDP & Error solución paper & Error solución notebook \\
\midrule
Adam   & 0.152 & 0.980 \\
Newton & 0.150 & 0.875 \\
\bottomrule
\end{tabular}
\label{tab:validacion-referencia}
\end{table}

Como no se ejecutó la escala completa del paper, esta comparación debe interpretarse con cuidado. El notebook reproduce el procedimiento general, pero al usar una red más pequeña y menos pasos de entrenamiento, la predicción queda más lejos de la referencia clásica.

\begin{figure}[H]
    \centering
    \includegraphics[width=1\linewidth]{figura5.png}
    \caption{error relativo de solución para PINNs}
    \label{fig:placeholder}
\end{figure}

Este mapa serviría para ver en qué zonas la red se equivoca más. En este problema, se esperaría que los errores aparezcan con más fuerza en regiones donde el campo cambia de forma brusca, especialmente cuando la viscosidad es baja.

\subsection{Sensibilidad}

El segundo criterio consiste en mirar si los resultados cambian cuando se modifica la configuración del experimento. En este informe, la sensibilidad se observa principalmente al comparar la escala usada en el paper con la escala usada en el notebook.\\

El paper usa una malla de $256\times256$, $72$ tiempos, una red de ancho $256$, profundidad $10$ y $100000$ pasos de entrenamiento. En cambio, el notebook en modo demo usa una malla de $48\times48$, $24$ tiempos, una red de ancho $80$, profundidad $4$ y $1500$ pasos. Esta diferencia explica por qué los errores del notebook son mucho mayores.\\

También se observa sensibilidad frente al número de datos usados. En el paper se usan $2048$, $8192$ y $32768$ datos para las PINNs, mientras que en el notebook demo solo se usan $512$ y $2048$. Al trabajar con menos datos y menos entrenamiento, la red tiene menos información para ajustar la solución y estimar los parámetros.

\begin{table}[H]
\centering
\scriptsize
\setlength{\tabcolsep}{3pt}
\caption{Sensibilidad de los errores frente al número de datos y al optimizador de la EDP en el notebook demo.}
\begin{tabular}{lcccc}
\toprule
Opt. EDP & $N_{data}$ & Error $\alpha$ & Error sol. & Error $\nu$ \\
\midrule
Adam   & 512  & 1.060 & 0.385 & 0.338 \\
Newton & 512  & 0.814 & 0.548 & 0.343 \\
Adam   & 2048 & 1.135 & 0.980 & 0.767 \\
Newton & 2048 & 7.300 & 0.875 & 0.553 \\
\bottomrule
\end{tabular}
\label{tab:sensibilidad-demo}
\end{table}

En los resultados obtenidos, Newton mejora un poco la estimación de $\nu$, pero no logra estabilizar bien la estimación de $\alpha$. Esto muestra que el comportamiento del método depende bastante de la configuración usada. En una escala reducida, el optimizador puede comportarse de manera distinta a como se comporta en el paper.

\subsection{Plausibilidad física}

El tercer criterio consiste en revisar si los resultados tienen sentido desde el punto de vista físico. Una solución plausible debe cumplir:

\begin{enumerate}[leftmargin=*]
\item las condiciones iniciales de las ecuaciones (7) y (8);
\item la anulación en la frontera del cuadrado;
\item residuos $f_u$ y $f_v$ pequeños en el dominio;
\item una evolución coherente con el balance entre convección y difusión;
\item mayor dificultad cuando $\nu$ disminuye y $\alpha/\nu$ aumenta.
\end{enumerate}

En este caso, la frontera se impone de forma exacta mediante el factor $x(1-x)y(1-y)$. Por eso, la red se anula automáticamente cuando $x=0$, $x=1$, $y=0$ o $y=1$. Esto ayuda a que la solución respete una parte importante del problema físico.

\begin{figure}[H]
    \centering
    \includegraphics[width=1\linewidth]{figura-perdida.png}
    \caption{Verificación física mediante la evolución de las pérdidas de entrenamiento. Se muestran la pérdida total, la pérdida asociada al residuo de la EDP, la pérdida de datos y la pérdida de condición inicial.}
    \label{fig:verificacion-fisica}
\end{figure}

La plausibilidad física también se observa en la dificultad del problema. Los errores tienden a ser mayores cuando $\nu$ es pequeño y $\alpha/\nu$ es grande (véase figuras 3, 4, 5, 6 y 7). Esto tiene sentido, porque en esos casos la difusión suaviza menos la solución y la convección domina más. Por eso aparecen cambios más bruscos en el campo de velocidades, lo cual hace que la red tenga más dificultad para aproximar la solución.\\

En conclusión, la validación muestra que el notebook reproduce la idea del paper, pero no alcanza su precisión. La comparación con FiPy muestra errores altos, la sensibilidad muestra que reducir la escala afecta bastante los resultados, y la plausibilidad física muestra que los errores aparecen justamente en los regímenes más difíciles.
\section{Limitaciones}

\begin{itemize}[leftmargin=*]
    \item \textbf{Escala reducida:} los resultados disponibles provienen de modo \texttt{demo}, no de la escala completa del paper, lo que provoca que el error aumente de una forma considerable.
    \item \textbf{Costo computacional:} ejecutar la grilla completa implica hasta 120 entrenamientos, lo cual puede ser pesado para Colab, por esto se decidió reducir la escala.
    \item \textbf{Tiempo de ejecución/costo temporal:} si se ejecuta la grilla completa, el tiempo de ejecución sería muy alto (puede tardar días/semanas), por lo que se debió limitar mucho la escala con la que se trabaja para poder tener un resultado dentro del tiempo estipulado, a pesar de que esto provocará un aumento del error.
    \item \textbf{Sensibilidad numerica:} los resultados de las PINNs dependen bastante de la configuración usada en el entrenamiento. Cambios en la arquitectura de la red, en los pesos de la función de pérdida, en la tasa de aprendizaje o en los puntos muestreados pueden modificar los errores obtenidos. Por esta razón, una misma formulación matemática puede producir resultados diferentes si se entrena con una red más pequeña, menos pasos o una distribución distinta de puntos. En este trabajo, esta limitación se observa al comparar la configuración completa del paper con el modo demo del notebook.

\end{itemize}

% ============================================================
\section{Conclusiones}

Se logró trabajar el problema de la ecuación de Burgers 2D usando redes neuronales informadas por la física. La idea principal fue estudiar cómo una red puede aprender el comportamiento de un fluido y, al mismo tiempo, intentar recuperar los parámetros que aparecen en la ecuación.\\

El notebook permitió seguir la metodología del paper. Se generaron datos de referencia, se entrenaron modelos y se compararon los resultados obtenidos. Sin embargo, no se hizo una réplica exacta del artículo, ya que se trabajó con una versión más pequeña para que el experimento pudiera ejecutarse en Colab.\\

La diferencia entre el paper y el notebook fue importante. El paper usa una configuración mucho más grande, mientras que el notebook usa menos puntos, una red más sencilla y menos tiempo de entrenamiento. Por esta razón, los errores obtenidos en el notebook fueron más altos.\\

La recuperación de los parámetros $\alpha$ y $\nu$ no fue precisa en la versión reducida. Esto muestra que, para este tipo de problemas, no basta con aplicar el método; también se necesita suficiente capacidad de cómputo para que la red pueda aprender mejor la solución.\\

También se observó que los casos con viscosidad más baja fueron más difíciles. Esto tiene sentido, porque cuando la viscosidad es pequeña el comportamiento del fluido cambia de forma más brusca y la red tiene más dificultad para aproximarlo correctamente.\\

La comparación entre la PINN y el método guiado solo por datos muestra una idea importante: parecerse a los datos no siempre significa entender bien la ecuación. Por eso, incluir la parte física dentro del entrenamiento sigue siendo una ventaja del enfoque PINN, aunque en este trabajo los resultados no hayan alcanzado la precisión del paper.\\

La extensión del trabajo permitió ver qué ocurre al reducir la escala del experimento. Se pudo ejecutar el procedimiento de forma más manejable, pero se perdió precisión. Esto muestra que el costo computacional influye bastante en los resultados.\\

Finalmente, el trabajo permitió entender mejor la propuesta del paper, sus ventajas y sus dificultades. Aunque los resultados del notebook no fueron iguales a los del artículo, sí sirvieron para reproducir la idea general, analizar las diferencias y reconocer las limitaciones de trabajar con una configuración reducida.


% ============================================================

\begin{thebibliography}{9}

\bibitem{new2024}
A. New, M. Villafane-Delgado y C. Shugert. \textit{Equation Identification for Fluid Flows via Physics-Informed Neural Networks}. arXiv:2408.17271, 2024.

\bibitem{raissi2019}
M. Raissi, P. Perdikaris y G. E. Karniadakis. \textit{Physics-informed neural networks: A deep learning framework for solving forward and inverse problems involving nonlinear partial differential equations}. Journal of Computational Physics, 378:686--707, 2019.

\bibitem{guyer2009}
J. E. Guyer, D. Wheeler y J. A. Warren. \textit{FiPy: Partial Differential Equations with Python}. Computing in Science \& Engineering, 11(3):6--15, 2009.

\bibitem{karniadakis2021}
G. E. Karniadakis, I. G. Kevrekidis, L. Lu, P. Perdikaris, S. Wang y L. Yang. \textit{Physics-informed machine learning}. Nature Reviews Physics, 3:422--440, 2021.

\end{thebibliography}

\end{document}
